\input{sol_pre.tex}
\input{q_2026-02-28_2.tex}
\begin{document}
\input{sol_head.tex}

%================= 1
\Savol{1}
$R=\underbrace{77\ldots7}_{2025}$ bölsin. Oxirgi raqami $d$ bölgan son $R+(d-7)$, \ $d=9,8,\dots,1$ (9 ta qöşiluvçi):
\[S=9R+\bigl((9+8+\dots+1)-9\cdot7\bigr)=9R-18.\]
$R=7\cdot\underbrace{11\ldots1}_{2025}$ 7 ga bölinadi $\Rightarrow$ $S\equiv-18\equiv3\pmod7$.
\Javob{D) 3}

%================= 2
\Savol{2}
$N=2025$, $M=2026$: \ $4\frac2N=\frac{4N+2}N$, \ $6\frac5M=\frac{6M+5}M$, \ $2\frac{2023}N=\frac{3N-2}N$, \ $7\frac{2021}M=\frac{8M-5}M$.
\[\frac{(4N+2)(6M+5)-(3N-2)(8M-5)}{NM}=\frac{35N+28M}{NM}=\frac{35}M+\frac{28}N.\]
$\frac{35}M+\frac{28}N-\frac{35}{2026}=\frac{28}{2025}$.
\Javob{D) $\frac{28}{2025}$}

%================= 3
\Savol{3}
Tezliklar $v$ va $v+2$: \ $\dfrac{24}v-\dfrac{24}{v+2}=1\Rightarrow48=v(v+2)\Rightarrow v=6$.\\
Ikkinçi yölovçi: $24:6=4$ soat.
\Javob{A) 4}

%================= 4
\Savol{4}
Tannarx $C$, narx $P$. 1-köylak: $0{,}95P-C=0{,}18\cdot0{,}95P\Rightarrow C=0{,}82\cdot0{,}95P=0{,}779P$.\\
2-köylak foydasi: $P-0{,}779P=0{,}221P$ — sotilgan narxning $22{,}1\%$.
\Javob{B) 22,1}

%================= 5
\Savol{5}
$(3^4\cdot2)^{-\frac12}=\dfrac1{9\sqrt2}$, \ $9^{-3}\cdot2^{-2}=\dfrac1{2916}$: \quad $\dfrac{2916}{9\sqrt2}=\dfrac{324}{\sqrt2}=162\sqrt2$.
\Javob{A) $162\sqrt2$}

%================= 6
\Savol{6}
$7-\sqrt{40}=(\sqrt5-\sqrt2)^2$ $\Rightarrow$ $\sqrt[6]{7-\sqrt{40}}=\sqrt[3]{\sqrt5-\sqrt2}$.
\[\sqrt[3]{(\sqrt5-\sqrt2)(\sqrt5+\sqrt2)}\cdot\sqrt[3]9=\sqrt[3]{3\cdot9}=3.\]
\Javob{A) 3}

%================= 7
\Savol{7}
$t=\sqrt[3]2$: \ $\sqrt[3]{32}=2t^2$, $\sqrt[3]{128}=4t$, \ $(t^2+1)(-t^2+2t+1)=5$ (çunki $t^3=2$, $t^4=2t$):
\[\frac5{2(t^2+1)}+\frac5{2(-t^2+2t+1)}=\frac{-t^2+2t+1}2+\frac{t^2+1}2=t+1;\qquad t+1-\sqrt[3]2=1.\]
\Javob{B) 1}

%================= 8
\Savol{8}
$a_1=3{,}1$, \ $a_{11}=3{,}8-7{,}7=-3{,}9$: \quad $S_{11}=\dfrac{3{,}1-3{,}9}2\cdot11=-4{,}4$.
\Javob{B) $-4{,}4$}

%================= 9
\Savol{9}
$b_1(1+q+q^2)=19$, \ $b_1^2(1+q^2+q^4)=133$; \ $1+q^2+q^4=(1+q+q^2)(1-q+q^2)$ $\Rightarrow$ $b_1(1-q+q^2)=7$.\\
Ayirsak: $2b_1q=12\Rightarrow b_2=6$; \ $b_1+b_3=13$, $b_1b_3=36$ $\Rightarrow$ kamayuvçi: $b_1=9$, $q=\frac23$.
\[S_5=\frac{9\left(1-\frac{32}{243}\right)}{1-\frac23}=27\cdot\frac{211}{243}=\frac{211}9=23\frac49.\]
\Javob{D) $23\frac49$}

%================= 10
\Savol{10}
$u=\sqrt a$, $v=\sqrt b$. Birinçi qavs (maxraj $(u-v)(u+v)$):
\[\frac{uv(u-v)+v^2(u+v)+2uv^2}{(u-v)(u+v)}=\frac{v(u+v)^2}{(u-v)(u+v)}=\frac{v(u+v)}{u-v}.\]
Ikkinçi qavs (maxraj $u(u-v)(u+v)$): \ $\dfrac{u(u+v)-v(u-v)-2uv}{u(u-v)(u+v)}=\dfrac{(u-v)^2}{u(u-v)(u+v)}=\dfrac{u-v}{u(u+v)}$.\\
Köpaytma: $\dfrac vu=\sqrt{\dfrac ba}=\sqrt{\dfrac{3\sqrt2}{\sqrt2}}=\sqrt3$.
\Javob{C) $\sqrt3$}

%================= 11
\Savol{11}
Bu Nyuton binomi: $\sum C_4^k(2x-1)^{4-k}(-2x)^k$ ($1,\,-4\cdot2x,\,6\cdot4x^2,\,-4\cdot8x^3,\,16x^4$):
\[\bigl((2x-1)-2x\bigr)^4=(-1)^4=1.\]
\Javob{A) 1}

%================= 12
\Savol{12}
$\cos\left(\frac{3\pi}2+2x\right)=\sin2x=0\Rightarrow x=\frac{\pi k}2$. \ $-\frac\pi2\le\frac{\pi k}2\le\frac{5\pi}2\Rightarrow k=-1,\dots,5$ — 7 ta.
\Javob{C) 7}

%================= 13
\Savol{13}
$\sin^4x+\cos^4x=1-\frac12\sin^22x=1\Rightarrow\sin2x=0\Rightarrow x=\frac{\pi k}2$.\\
$-\frac{3\pi}2\le\frac{\pi k}2\le\frac{5\pi}2\Rightarrow k=-3,\dots,5$ — 9 ta.
\Javob{C) 9}

%================= 14
\Savol{14}
$x>0$, $x\ne1$ (aks holda $1=9$). Logarifmlab bölamiz: $\dfrac{3y-4}{y+1}=\dfrac{\log_x9}{\log_x27}=\dfrac23\Rightarrow9y-12=2y+2\Rightarrow y=2$.\\
$x^3=27\Rightarrow x=3$. \ $xy=6$.
\Javob{D) 6}

%================= 15
\Savol{15}
Maxrajga köpaytirib: $-\sqrt{5x^3-4x^2-4x}=\log_5^2(5x^3-4x^2+1)$. Çap $\le0$, öng $\ge0$ $\Rightarrow$ ikkalasi 0:
\[5x^3-4x^2-4x=0,\qquad5x^3-4x^2+1=1\ \Rightarrow\ 4x=0,\ x=0.\]
Tekşiriş: maxraj $\sqrt4+0=2\ne0$. \ 1 ta ildiz.
\Javob{B) 1}

%================= 16
\Savol{16}
$\dfrac{2027!}{2025!(2026+1)}=\dfrac{2027\cdot2026}{2027}=2026$: \ $x^4-2026x^2+2025=0\Rightarrow x^2=1$ yoki $x^2=2025$.\\
$x=\pm1,\ \pm45$ — 4 ta.
\Javob{D) 4 ta}

%================= 17
\Savol{17}
$1+2+\dots+k=\frac{k(k+1)}2$; \ $x$ ni ham $\frac x1$ ($k=1$) deb olsak:
\[x\sum_{k=1}^{4041}\frac2{k(k+1)}=2x\left(1-\frac1{4042}\right)=\frac{4041}{2021}x=4041\ \Rightarrow\ x=2021.\]
\Javob{A) 2021}

%================= 18
\Savol{18}
Surat $>0$: $x\ne2$, $x\ne6$. \ Maxraj $>0$: $x^2-7x-18<0\Rightarrow-2<x<9$.\\
Butunlar: $-1,0,1,3,4,5,7,8$; yiğindi $35-2-6=27$.
\Javob{B) 27}

%================= 19
\Savol{19}
$x\le2$, \ $\sqrt{(x-2)^2}=2-x$. \ $t=\sqrt{2-x}\ge0$: \ $t^2<5t+6\Rightarrow(t-6)(t+1)<0\Rightarrow t<6$.\\
$2-x<36\Rightarrow x>-34$. \ $-33\le x\le2$ — 36 ta.
\Javob{B) 36}

%================= 20
\Savol{20}
$x+2>0$, $x+2\ne1$, $x\ne2$, \ $\sqrt{24+2x-x^2}>0\Rightarrow-4<x<6$.\\
$x\in(-2;6)$, $x\ne-1$, $x\ne2$: \ $0,1,3,4,5$ — 5 ta.
\Javob{A) 5 ta}

%================= 21
\Savol{21}
$\dfrac{\sin^3x}{\cos^7x}$ — toq funksiya, oraliq simmetrik $\Rightarrow$ integral 0.
\Javob{B) 0}

%================= 22
\Savol{22}
$t=x^2+1\ge1$: \ $x^4+1=(t-1)^2+1$, \ $y=t-2+\dfrac2t\ge2\sqrt2-2$ ($t=\sqrt2$ da tenglik).\\
$t\to\infty$ da $y\to\infty$. \ $E(y)=[2\sqrt2-2;\infty)$.
\Javob{C) $[2\sqrt2-2;\infty)$}

%================= 23
\Savol{23}
Markaziy burçak $90^\circ$ $\Rightarrow$ $\triangle AOB$ teng yonli töğri burçakli: $AB=R\sqrt2=8\Rightarrow R=4\sqrt2$.
\Javob{C) $4\sqrt2$}

%================= 24
\Savol{24}
$f'(x)=e^{x-7}(x-7)=0\Rightarrow x=7$. \ $f(5)=-\frac3{e^2}$, $f(7)=-1$, $f(9)=e^2$. Eng kiçigi $-1$.
\Javob{A) $-1$}

%================= 25
\Savol{25}
$AF=6$, $FC=12$; $CD=6$, $DB=12$; $BE=6$, $EA=12$. Burçaklardagi uçta uçburçak teng:
\[S_{CFD}=\frac12\cdot12\cdot6\sin60^\circ=18\sqrt3,\qquad S_{ABC}=\frac{\sqrt3}4\cdot18^2=81\sqrt3.\]
$S_{FED}=81\sqrt3-3\cdot18\sqrt3=27\sqrt3$.
\Javob{A) $27\sqrt3$}

%================= 26
\Savol{26}
$\alpha,\alpha$ — bissektrisa: $\frac{AB}{AC}=\frac{6{,}5}{7{,}5}\Rightarrow AB=13$. Uçburçak $13,14,15$: \ $p=21$, $S=\sqrt{21\cdot8\cdot7\cdot6}=84$.\\
Ikki bissektrisa $I$ — içki aylana markazida kesişadi, $r=\frac Sp=4$.
\fig{\begin{tikzpicture}[scale=0.3]
\coordinate (B) at (0,0); \coordinate (C) at (14,0); \coordinate (A) at (5,12); \coordinate (I) at (6,4); \coordinate (D) at (6.5,0); \coordinate (K) at (9.2,6.4);
\fill[okgreen!30] (A)--(I)--(C)--cycle;
\draw (A)--(B)--(C)--cycle; \draw (A)--(D); \draw (C)--(2.414,5.793);
\draw[axs] (I)--(K); \draw[ax] (I) circle (4);
\fill (I) circle (4pt);
\node[above] at (A){$A$}; \node[below left] at (B){$B$}; \node[below right] at (C){$C$}; \node[below] at (D){$D$}; \node[below] at (I){$I$};
\node[right] at (9.6,6.8){$15$}; \node[below,font=\scriptsize] at (3.25,0){$6{,}5$}; \node[below,font=\scriptsize] at (10.25,0){$7{,}5$};
\node[aux,font=\scriptsize] at (7.2,5.8){$r$};
\end{tikzpicture}}
Böyalgan soha $=S_{AIC}=\frac12\cdot AC\cdot r=\frac12\cdot15\cdot4=30$.
\Javob{D) 30}

%================= 27
\Savol{27}
$BC=3$, $AD=9$, $AB=6$ — balandlik. $\triangle BOC\sim\triangle DOA$, $k=\frac39=\frac13$ $\Rightarrow$ ularning balandliklari $1:3$.
\fig{\begin{tikzpicture}[scale=0.42]
\coordinate (A) at (0,0); \coordinate (D) at (9,0); \coordinate (B) at (0,6); \coordinate (C) at (3,6); \coordinate (O) at (2.25,4.5);
\draw (A)--(B)--(C)--(D)--cycle; \draw (A)--(C); \draw (B)--(D);
\draw[axs] (O)--(2.25,0); \draw[ax] (O)--(2.25,6); \draw (0,0.4)--(0.4,0.4)--(0.4,0);
\node[below left] at (A){$A$}; \node[above left] at (B){$B$}; \node[above] at (C){$C$}; \node[below right] at (D){$D$}; \node[right] at (O){$O$};
\node[left] at (0,3){$6$}; \node[below] at (5,0){$9$};
\end{tikzpicture}}
$O$ dan $AD$ gaça: $6\cdot\dfrac34=4{,}5$.
\Javob{A) 4,5}

%================= 28
\Savol{28}
$\dfrac{n(n-3)}2=35\Rightarrow n=10$. Taşqi aylanada har tomon $36^\circ$ yoyni tortadi. Eng kiçik diagonal 2-uçga, eng kattasi (diametr) 5-uçga boradi; ular orasidagi içki çizilgan burçak $3\cdot36^\circ$ yoyga tiraladi: \ $\frac{108^\circ}2=54^\circ$.
\Javob{A) $54^\circ$}

%================= 29
\Savol{29}
$AB$ örtasi $M(1;1;3)$, \ $\overrightarrow{CM}=(1;-1;1)$, \ $CM=\sqrt3$.
\Javob{D) $\sqrt3$}

%================= 30
\Savol{30}
Yon qirra $l=\sqrt{39}$, $60^\circ$: \ $H=l\sin60^\circ=\frac{3\sqrt{13}}2$, \ asos taşqi radiusi $R=l\cos60^\circ=\frac{\sqrt{39}}2$ — oltiburçak tomoni.
\[S=\frac{3\sqrt3}2\cdot\frac{39}4=\frac{117\sqrt3}8,\qquad V=\frac13\cdot\frac{117\sqrt3}8\cdot\frac{3\sqrt{13}}2=\frac{117\sqrt{39}}{16}.\]
\Javob{D) $\frac{117\sqrt{39}}{16}$}

%================= 31
\Savol{31}
Kerakli element $A$ yoki $C$ ning birida va $B$ yoki $D$ ning birida bölişi kerak. Rasmdan:\\
$A\cap B$: $19,14$; \ $C\cap B$: $1,6,9$; \ $A\cap B\cap C$: $24$; \ $A\cap D$: $5,25,10$; \ $A\cap C\cap D$: $18$; \ $C\cap D$: $8,13$.\\
Jami: $2+3+1+3+1+2=12$.
\Javob{A) 12}

%================= 32
\Savol{32}
10 harfdan 2 tasi M: \ $P=\frac2{10}=\frac15$.
\Javob{C) $\frac15$}

%================= 33-35
\newpage
\Savol{33}
Konus asosi — asos uçburçakning içki aylanasi ($r=4$), uçi — piramida uçi. $SK$ — piramida apofemasi, u konus yasovçisi ham.
\fig{\begin{tikzpicture}[scale=0.45]
\coordinate (O) at (0,0); \coordinate (K) at (4,0); \coordinate (S) at (0,6.93);
\draw (S)--(K)--(O); \draw[axs] (S)--(O); \draw (0,0.4)--(0.4,0.4)--(0.4,0);
\draw (3.4,0) arc (180:120:0.6); \node[font=\scriptsize] at (3.0,0.5){$60^\circ$};
\node[below left] at (O){$O$}; \node[below right] at (K){$K$}; \node[above] at (S){$S$}; \node[below] at (2,0){$4$};
\end{tikzpicture}}
$SO=4\tg60^\circ=4\sqrt3$, \ $SK=\frac4{\cos60^\circ}=8$. \ Asos tomoni $a=2\sqrt3r=8\sqrt3$, \ $S_{asos}=\frac{\sqrt3}4\cdot192=48\sqrt3$.\\
\textbf{33.} $S_{kon}=\pi rl=\pi\cdot4\cdot8=32\pi\approx96$.\\
\textbf{34.} $S_{yon}=\dfrac{S_{asos}}{\cos60^\circ}=96\sqrt3$ (yoki $3\cdot\frac12\cdot8\sqrt3\cdot8$).\\
\textbf{35.} $V=\frac13\cdot48\sqrt3\cdot4\sqrt3=192$.
\Javob{33 — B; \quad 34 — E; \quad 35 — C}

%================= 36
\newpage
\Savol{36}
Ikkinçi tenglamani 2 ga köpaytiramiz: $2x^2-5xy+2y^2=0\iff(2x-y)(x-2y)=0$.

1) $y=2x$: \ $x-\frac12x^2-2x=0\Rightarrow x\left(1+\frac x2\right)=0$ $\Rightarrow$ $(0;0)$, $(-2;-4)$.

2) $x=2y$: \ $2y-\frac12y^2-y=0\Rightarrow y\left(1-\frac y2\right)=0$ $\Rightarrow$ $(0;0)$, $(4;2)$.

Yeçimlar: $(0;0)$, $(-2;-4)$, $(4;2)$.

a) $xy$: $0,\ 8,\ 8$ $\Rightarrow$ eng kattasi 8. \qquad b) 3 ta yeçim.
\Javob{a) 8; \quad b) 3}

%================= 37
\newpage
\Savol{37}
$\cos x=0$ bölsa, $\sin^2x=1\ne0$ — ildiz emas. $\cos^2x$ ga bölamiz:
\[\tg^2x+2\tg x-3=0\ \Rightarrow\ \tg x=1\ \text{yoki}\ \tg x=-3;\qquad x=\frac\pi4+\pi k,\quad x=-\arctg3+\pi k.\]
a) Manfiylari: $-\frac{3\pi}4\approx-2{,}36$, \ $-\arctg3\approx-1{,}25$ $\Rightarrow$ eng kattasi $-\arctg3$.

b) $[-\pi;\pi]$ da: $-\frac{3\pi}4,\ \frac\pi4,\ -\arctg3,\ \pi-\arctg3$ — 4 ta.
\Javob{a) $-\arctg3$; \quad b) 4}

%================= 38
\newpage
\Savol{38}
Ikkinçi tenglikda $x$ örniga $3x-3$ qöyamiz: $x+1\to3x-2$, \ $\frac x3-4\to x-5$:
\[f(3x-2)-g(x-5)=9x-9.\]
Birinçi tenglikdan ayiramiz: $8g(x-5)=x+1-9x+9\Rightarrow g(x-5)=\frac54-x$, ya'ni $g(t)=-t-\frac{15}4$.\\
$f(3x-2)=9x-9+\frac54-x=8x-\frac{31}4$; \ $t=3x-2$: \ $f(t)=\frac{8(t+2)}3-\frac{31}4=\frac{8t}3-\frac{29}{12}$.

a) $f(1)=\frac83-\frac{29}{12}=\frac3{12}=\frac14$ $\Rightarrow$ $(f(1))^{-1}=4$.

b) $f(x)+g(x)=\frac{8x}3-x-\frac{29}{12}-\frac{45}{12}=\frac{5x}3-\frac{37}6=\frac12$ — çiziqli tenglama, 1 ta ildiz ($x=4$).
\Javob{a) 4; \quad b) 1}

%================= 39
\newpage
\Savol{39}
Grafikdan (katak — 1 birlik): lokal ekstremumlar $x=-8$ (min), $-7$ (max), $-6$ (min), $-5$ (max), $-3$ (min), $0$ (max), $2$ (min), $4$ (max).

a) $[-9;5]$ kesmaning içki nuqtalarida 8 ta ekstremum.

b) $f'(x)<0$ — funksiya kamayadigan oraliqlar: $[-9;-8)$, $(-7;-6)$, $(-5;-3)$, $(0;2)$, $(4;5]$.\\
Ekstremum nuqtalarda $f'=0$. Butun sonlar: $-9$, $-4$, $1$, $5$ — 4 ta.
\Javob{a) 8; \quad b) 4}

%================= 40
\newpage
\Savol{40}
$g(x)$ nollari: $0,\ -\frac53,\ 2$ $\Rightarrow$ $g(x)=ax\left(x+\frac53\right)(x-2)$; \ $g(-2)=a(-2)\left(-\frac13\right)(-4)=-\frac83a=-8\Rightarrow a=3$.
\[g(x)=x(3x+5)(x-2)=3x^3-x^2-10x.\]
$f(2)=4p+2q=0$, \ $f(-2)=4p-2q=-8$ $\Rightarrow$ $p=-1$, $q=2$, \ $f(x)=-x^2+2x$.

a) $\dfrac{a+b+c}{p+q}=\dfrac{3-1-10}{-1+2}=-8$.

b) $g(x)-f(x)=3x^3-12x=3x(x-2)(x+2)$: $[-2;0]$ da $\ge0$, $[0;2]$ da $\le0$.
\[S=\int_{-2}^0(3x^3-12x)\,dx+\int_0^2(12x-3x^3)\,dx=\left(0-(12-24)\right)+(24-12)=12+12=24.\]
\Javob{a) $-8$; \quad b) 24}

%================= 41
\newpage
\Savol{41}
$CO$ — $C$ töğri burçak bissektrisasi. $O$ dan $AB$ ga perpendikulyar $OP=r=\sqrt3$.
\fig{\begin{tikzpicture}[scale=0.45]
\coordinate (C) at (0,0); \coordinate (B) at (14.888,0); \coordinate (A) at (0,3.989); \coordinate (O) at (1.732,1.732); \coordinate (D) at (3.146,3.146); \coordinate (P) at (2.180,3.405);
\draw (A)--(B)--(C)--cycle; \draw (O) circle (1.732); \draw (C)--(D);
\draw[axs] (O)--(P); \draw[ax] (O)--(1.732,0);
\draw (0,0.45)--(0.45,0.45)--(0.45,0); \fill (O) circle (3pt);
\node[left] at (A){$A$}; \node[below left] at (C){$C$}; \node[below right] at (B){$B$}; \node[above right] at (D){$D$}; \node[below right] at (O){$O$}; \node[above left] at (P){$P$};
\end{tikzpicture}}
a) $CO=r\sqrt2=\sqrt6$ (kvadrat diagonali), \ $OD=CO\cdot\dfrac{\sqrt2}{\sqrt3}=\sqrt6\cdot\dfrac{\sqrt2}{\sqrt3}=2$.

b) $\triangle OPD$: $\sin\angle ODP=\dfrac{OP}{OD}=\dfrac{\sqrt3}2$ $\Rightarrow$ $CD$ va $AB$ orasidagi burçak $60^\circ$ (qöşnisi $120^\circ$).\\
$\triangle CDB$: $\angle DCB=45^\circ$; \ $\angle CDB=120^\circ$ bölsa $\angle B=15^\circ$, \ $\angle CDB=60^\circ$ bölsa $\angle B=75^\circ$, $\angle A=15^\circ$.\\
Har ikki holda ötkir burçaklar $15^\circ$ va $75^\circ$; kiçigi $15^\circ$.
\Javob{a) 2; \quad b) $15^\circ$}

%================= 42
\newpage
\Savol{42}
Kvadratlar umumiy uçi $P$; kiçik kvadrat $PRQT$ (tomoni $a$), $Q$ — katta kvadratning yuqori tomonida. $PT$ va katta kvadratning çap tomoni orasidagi burçak $\theta$, $M$ — $TQ$ va çap tomon kesişişi.
\fig{\begin{tikzpicture}[scale=0.36]
\coordinate (P) at (0,0); \coordinate (R) at (5.367,2.683); \coordinate (T) at (-2.683,5.367); \coordinate (Q) at (2.683,8.050); \coordinate (M) at (0,6.708);
\draw (0,0) rectangle (8.05,8.05); \draw (P)--(R)--(Q)--(T)--cycle;
\draw[aux,thick] (0,1.2) arc (90:116.57:1.2); \node[aux,font=\scriptsize] at (-0.35,1.6){$\theta$};
\draw[ax] (Q)--(2.683,0);
\node[below] at (P){$P$}; \node[right] at (R){$R$}; \node[above] at (Q){$Q$}; \node[left] at (T){$T$}; \node[right,font=\scriptsize] at (M){$M$};
\node at (-0.9,4.2){$S_1$}; \node at (1.6,4.5){$S_2$}; \node[font=\scriptsize] at (0.6,7.6){$S_3$}; \node at (6,6){$S_4$};
\end{tikzpicture}}
$\triangle PTM$: $\angle T=90^\circ$, $PT=a$ $\Rightarrow$ $TM=a\tg\theta$, \ $S_1=\frac12a^2\tg\theta$. \ $S_1+S_2=a^2$ $\Rightarrow$ $\dfrac{S_1}{a^2}=\dfrac9{36}$:
\[\frac12\tg\theta=\frac14\ \Rightarrow\ \tg\theta=\frac12,\quad a^2=36,\ a=6,\quad \cos\theta=\frac2{\sqrt5},\ \sin\theta=\frac1{\sqrt5}.\]
Katta kvadrat tomoni — $Q$ ning balandligi: $L=a\cos\theta+a\sin\theta=\frac{18}{\sqrt5}$, \ $L^2=\frac{324}5$; \ $PM=\frac a{\cos\theta}=3\sqrt5$.

a) $S_3$ — katetlari $L-PM=\frac{18}{\sqrt5}-3\sqrt5=\frac3{\sqrt5}$ va $Q$ gaça gorizontal masofa $a\cos\theta-a\sin\theta=\frac6{\sqrt5}$ bölgan töğri burçakli uçburçak:
\[S_3=\frac12\cdot\frac3{\sqrt5}\cdot\frac6{\sqrt5}=\frac95.\]
b) $S_4=L^2-S_2-S_3=\frac{324}5-27-\frac95=36$.
\Javob{a) $\frac95$; \quad b) 36}

%================= 43
\newpage
\Savol{43}
Muntazam oltiburçak içki burçagi $120^\circ$. $BG=1$, $AB=9$, $AH=9-2=7$. \ $AG$ ni ötkazamiz.
\fig{\begin{tikzpicture}[scale=0.25]
\coordinate (A) at (0,0); \coordinate (F) at (9,0); \coordinate (E) at (13.5,7.794); \coordinate (D) at (9,15.588); \coordinate (C) at (0,15.588); \coordinate (B) at (-4.5,7.794);
\coordinate (G) at (-4,8.66); \coordinate (H) at (7,0);
\fill[okgreen!30] (G)--(B)--(A)--(H)--cycle;
\draw (A)--(B)--(C)--(D)--(E)--(F)--cycle; \draw (G)--(H); \draw[axs] (G)--(A);
\node[below left] at (A){$A$}; \node[left] at (B){$B$}; \node[above left] at (C){$C$}; \node[above right] at (D){$D$}; \node[right] at (E){$E$}; \node[below right] at (F){$F$};
\node[above left] at (G){$G$}; \node[below] at (H){$H$}; \node at (-0.5,4.2){$S_1$}; \node at (6,9){$S_2$};
\end{tikzpicture}}
$\triangle GBA$: $AG^2=1+81-2\cdot9\cos120^\circ=91$; \ $\varphi=\angle GAB$: \ $\cos\varphi=\dfrac{91+81-1}{2\cdot9\sqrt{91}}=\dfrac{19}{2\sqrt{91}}$, \ $\sin\varphi=\dfrac{\sqrt3}{2\sqrt{91}}$.\\
$\angle GAH=120^\circ-\varphi$: \ $\cos\angle GAH=-\frac12\cdot\frac{19}{2\sqrt{91}}+\frac{\sqrt3}2\cdot\frac{\sqrt3}{2\sqrt{91}}=-\frac4{\sqrt{91}}$, \ $\sin\angle GAH=\frac{5\sqrt3}{\sqrt{91}}$.

a) $GH^2=91+49+2\cdot\sqrt{91}\cdot7\cdot\frac4{\sqrt{91}}=196\Rightarrow GH=14$.

b) $S_1=S_{GBA}+S_{GAH}=\frac12\cdot1\cdot9\cdot\frac{\sqrt3}2+\frac12\cdot\sqrt{91}\cdot7\cdot\frac{5\sqrt3}{\sqrt{91}}=\frac{9\sqrt3}4+\frac{70\sqrt3}4=\frac{79\sqrt3}4$.\\
$S_{6}=\frac{3\sqrt3}2\cdot81=\frac{486\sqrt3}4$, \ $S_2=\frac{407\sqrt3}4$ $\Rightarrow$ $\dfrac{S_1}{S_2}=\dfrac{79}{407}$.
\Javob{a) 14; \quad b) $\frac{79}{407}$}

%================= 44
\newpage
\Savol{44}
$CH$ — $AB$ ga mediana ($AC=BC$ $\Rightarrow$ $CH\perp AB$), $M$ — medianalar kesişgan nuqta, $S$ — sfera markazi, $SM\perp ABC$, $SM=\frac23$.
\fig{\begin{tikzpicture}[scale=1.1]
\coordinate (A) at (0,0); \coordinate (B) at (4,0); \coordinate (C) at (2.6,1.6); \coordinate (D) at (0,3.2);
\coordinate (H) at (2,0); \coordinate (M) at (2.2,0.533); \coordinate (S) at (2.2,1.13);
\draw (D)--(A)--(B)--(D)--(C)--(B); \draw[dashed] (A)--(C);
\draw[okgreen,thick] (S) circle (0.6);
\draw[ax] (C)--(H); \draw[ax] (S)--(M); \draw[ax] (D)--(H);
\fill (M) circle (1pt); \fill (S) circle (1pt);
\node[below left] at (A){$A$}; \node[below right] at (B){$B$}; \node[right] at (C){$C$}; \node[above] at (D){$D$}; \node[below] at (H){$H$}; \node[below right,font=\scriptsize] at (M){$M$};
\end{tikzpicture}}
a) $AD\perp ABC$ $\Rightarrow$ $ABD\perp ABC$; $MH\perp AB$, şuning uçun $S$ dan $ABD$ gaça masofa $=MH=\frac23$. \ $MH=\frac13CH\Rightarrow CH=2$.

b) $C$ dan sferaga urinmalar teng: $CK=CM=\frac23CH=\frac43$ ($K$ — $DC$ dagi uriniş nuqtasi) $\Rightarrow$ $DK=3-\frac43=\frac53$.\\
Sfera $ABD$ ga $T$ da urinadi: $T$ — $H$ ustida, $AB$ dan $\frac23$ balandlikda. $D$ dan urinmalar teng: $DT=DK=\frac53$.\\
$AD=h$, $AH=x$: \ $DC^2=h^2+x^2+CH^2\Rightarrow h^2+x^2=5$; \ $DT^2=x^2+\left(h-\frac23\right)^2=\frac{25}9$.
\[5-\frac43h+\frac49=\frac{25}9\ \Rightarrow\ h=2,\quad x=1,\ AB=2.\]
$V=\frac13\cdot\frac12\cdot AB\cdot CH\cdot AD=\frac13\cdot2\cdot2=\frac43$.
\Javob{a) 2; \quad b) $\frac43$}

%================= 45
\newpage
\Savol{45}
Foyda: $A(x)-B(x)=\frac{2x}3-2\sqrt x=\frac23\sqrt x\left(\sqrt x-3\right)$.

a) Zarar körmaslik: $A(x)\ge B(x)\iff\sqrt x\ge3\iff x\ge9$ $\Rightarrow$ kamida 9 işçi ($x=9$ da foyda 0).

b) Foyda: $A(x)>B(x)\iff x>9$ $\Rightarrow$ eng kamida 10 işçi.
\Javob{a) 9; \quad b) 10}

\ShowJavoblar
\end{document}
